\chapter{Worked example: pseudorandom functions}

Following Rosulek, \emph{The Joy of Cryptography}, §6. A pseudorandom function is
\emph{perfectly} secure when it is built from a family of bijections: for each
input the keyed map is a bijection of the key space onto the output space, so a
uniform key produces a uniform output and the real (keyed) and ideal (random)
oracles are indistinguishable with advantage exactly $0$. As with the one-time pad
this is a single bijection coupling, over an abstract, possibly cross-type,
key/output pair.

\begin{definition}[Bijection PRF family]
  \label{def:bijprf}
  \uses{def:spcomp}
  \lean{CatCrypt.Examples.PRF.BijPRFFamily, CatCrypt.Examples.PRF.BijPRFFamily.toPRFScheme}
  \leanok
  A bijection PRF family has finite nonempty key, input, and output types together
  with, for each input $x$, a bijection $\mathsf{bij}_x : \mathrm{Key} \simeq
  \mathrm{Output}$. It induces a PRF scheme evaluating $\mathsf{eval}(k, x) =
  \mathsf{bij}_x(k)$.
\end{definition}

\begin{lemma}[Real and ideal oracles couple]
  \label{lem:bijprf-coupling}
  \uses{def:bijprf}
  \lean{CatCrypt.Examples.PRF.bijPRF_coupling}
  \leanok
  For every input $x$, the real oracle (a uniform key run through
  $\mathsf{bij}_x$) and the ideal oracle (a uniform output) are related by the
  bijection coupling $\mathsf{bij}_x$.
\end{lemma}

\begin{theorem}[Perfect PRF security]
  \label{thm:bijprf-perfect}
  \uses{lem:bijprf-coupling}
  \lean{CatCrypt.Examples.PRF.bijPRF_perfect}
  \leanok
  A bijection PRF family has perfect security: for every input $x$ and adversary
  $A$, the PRF advantage $\mathsf{Adv}^{\mathrm{PRF}}(x, A) = 0$.
\end{theorem}

\begin{corollary}[XOR PRF]
  \label{cor:xorprf}
  \uses{thm:bijprf-perfect}
  \lean{CatCrypt.Examples.PRF.boolXorPRF, CatCrypt.Examples.PRF.boolXorPRF_perfect}
  \leanok
  The XOR pseudorandom function $F(k, x) = k \oplus x$, instantiating the family by
  $x \mapsto (k \mapsto k \oplus x)$, has perfect PRF security.
\end{corollary}

\begin{theorem}[Cascade PRF security bound]
  \label{thm:cascade-prf}
  \uses{def:bijprf}
  \lean{CatCrypt.Examples.PRF.cascade_prf_bound}
  \leanok
  The \emph{computational} cascade result. For the double construction
  $\mathrm{cascade}(k_1, k_2, x) = F(k_2, F(k_1, x))$, composing the two
  PRF-security hops and the PRF/PRP switching hop through the advantage triangle
  inequality gives
  \[
    \mathsf{Adv}^{\mathrm{PRF}}_{\mathrm{cascade}} \;\le\; 2\,\varepsilon_{\mathrm{prf}}
    + \frac{q(q-1)}{2N},
  \]
  where the birthday term is \emph{proved} (via the switching lemma), not assumed.
  The three hops are the named cryptographic hypotheses, discharged by a keyed
  $\varepsilon$-PRF-security assumption.
\end{theorem}
